% Zone „Das kennst du schon“ – aus bank/bruchrechnung/zone.jsonl (zusammenbau v0.1)
\einheitenkopf[zone][Kennst du schon]{Das kennst du schon}
\setcounter{aufgabe}{0}

%% TODO Titel: Ich-kann-Satz fehlt in der Bank (Platzhalter: Kettenname) – Nr. 1
%% TODO Anweisung: ein Satz über der ersten Teilaufgabe fehlt in der Bank – Nr. 1
\begin{aufgabe}{Bruch als Anteil lesen, am Streifen einzeichnen (drei Achtel markieren)}
\begin{teile}
\teil Welcher Anteil ist gefärbt?

\bruchrechteck{2}{5}
\leerfeld
\teil Der Streifen hat 12 Teile. Färbe $\frac{2}{3}$ davon.

\bruchrechteck{0}{12}
\end{teile}
\end{aufgabe}

%% TODO Titel: Ich-kann-Satz fehlt in der Bank (Platzhalter: Kettenname) – Nr. 2
%% TODO Anweisung: ein Satz über der ersten Teilaufgabe fehlt in der Bank – Nr. 2
\begin{aufgabe}{Kürzen und Erweitern; gleichwertige Brüche (zwei Drittel gleich vier Sechstel)}
\begin{teile}
\teil Erweitere $\frac{2}{5}$ mit 3.
\teil Kürze $\frac{24}{36}$ so weit wie möglich.
\end{teile}
\end{aufgabe}

%% TODO Titel: Ich-kann-Satz fehlt in der Bank (Platzhalter: Kettenname) – Nr. 3
%% TODO Anweisung: ein Satz über der ersten Teilaufgabe fehlt in der Bank – Nr. 3
\begin{aufgabe}{Vielfache und Teiler (kgV von vier und sechs, ggT von zwölf und achtzehn)}
\begin{teile}
\teil Nenne die ersten vier Vielfachen von 3.
\teil Größter gemeinsamer Teiler von 16 und 24?
\end{teile}
\end{aufgabe}

%% TODO Titel: Ich-kann-Satz fehlt in der Bank (Platzhalter: Kettenname) – Nr. 4
%% TODO Anweisung: ein Satz über der ersten Teilaufgabe fehlt in der Bank – Nr. 4
\begin{aufgabe}{Stellenwerte der Dezimalzahlen (Zehntel, Hundertstel; 0,7 = 0,70)}
\begin{teile}
\teil 4 Zehntel als Dezimalzahl
\teil 2 Einer, 5 Zehntel und 8 Hundertstel als Dezimalzahl
\end{teile}
\end{aufgabe}

%% TODO Titel: Ich-kann-Satz fehlt in der Bank (Platzhalter: Kettenname) – Nr. 5
%% TODO Anweisung: ein Satz über der ersten Teilaufgabe fehlt in der Bank – Nr. 5
\begin{aufgabe}{Schriftliches Rechnen mit natürlichen Zahlen, Einmaleins}
\begin{teile}
\teil $7 \cdot 8$
\teil Schriftlich: $3\,475\text{ g} + 1\,268\text{ g}$
\end{teile}
\end{aufgabe}

%% TODO Titel: Ich-kann-Satz fehlt in der Bank (Platzhalter: Kettenname) – Nr. 6
%% TODO Anweisung: ein Satz über der ersten Teilaufgabe fehlt in der Bank – Nr. 6
\begin{aufgabe}{Bruch ↔ Dezimalzahl bei einfachen Brüchen (ein Halb, drei Viertel)}
\begin{teile}
\teil $\frac{3}{10}$ als Dezimalzahl
\teil $\frac{3}{5}$ als Dezimalzahl
\end{teile}
\end{aufgabe}

%% TODO Titel: Ich-kann-Satz fehlt in der Bank (Platzhalter: Kettenname) – Nr. 7
%% TODO Anweisung: ein Satz über der ersten Teilaufgabe fehlt in der Bank – Nr. 7
\begin{aufgabe}{Kürzen und Erweitern; gleichwertige Brüche (zwei Drittel gleich vier Sechstel) – Fehler finden}
\begin{teile}
\teil Lea erweitert: $\frac{3}{5} = \frac{3}{10}$. Finde den Fehler.
\end{teile}
\end{aufgabe}

%% TODO Titel: Ich-kann-Satz fehlt in der Bank (Platzhalter: Kettenname) – Nr. 8
%% TODO Anweisung: ein Satz über der ersten Teilaufgabe fehlt in der Bank – Nr. 8
\begin{aufgabe}{Kürzen und Erweitern; gleichwertige Brüche (zwei Drittel gleich vier Sechstel) – selbst rechnen}
\begin{teile}
\teil Schreibe $\frac{4}{9}$ mit dem Nenner 27.
\end{teile}
\end{aufgabe}
